\section{Conclusion}

The climate of a world with Earth-like air and water depends on how familiar processes interact under unfamiliar
planetary conditions. Adapting ExoPlaSim, CM1, SWAN and SWASH to a breathable Moon required following those
interactions from the global circulation into clouds, electrified storms and waves at the shore. In doing so,
the work connected two uses of planetary models: investigating a world still being designed and examining the
assumptions used to compute the climates of natural planets.

That connection arose from the combination of conditions in the design. Weak gravity puts a large mass of air
above each square metre at a familiar surface pressure, and raises the pressure levels through a substantial
fraction of a small world's radius; a long solar day gives clouds, soil and storms a different duration of
heating. Several approximations made for Earth were therefore strained together. The same runs also made
ExoPlaSim's clock, rain-unit and orbit defects conspicuous, because their errors grew with the long day and small
tilt. By the code's arithmetic those defects also affect natural-planet configurations that meet their
conditions. A constructed world can therefore serve as an experiment on the tools used to study worlds found
around other stars.

The Sun-clock and orbit tests illustrate why the consequence extends beyond the accuracy of a mean climate.
On the test world those defects redistributed sunlight while preserving its global mean at the top of the air.
Correcting them changed the global-mean surface temperature by less than a kelvin over the two-year test, while
the fraction of land receiving little rain grew fourfold. A small change in mean warmth accompanied a large
change in where water reached the land. For a natural planet, that distinction bears on the distribution of
potentially habitable environments; for a constructed one, it bears on which basins could sustain the lakes
included in its design. Comparing the released and repaired code on slowly rotating planets would measure how
far this consequence extends beyond the lunar case.

The larger changes in the test world's mean temperature came from the empirical radiation fits and cloud
settings. Their sensitivity makes the meaning of Earth-like air especially consequential: a gas mixture can
resemble Earth's while its absorber paths lie beyond those used to fit a radiation scheme, and water can
condense under familiar temperatures while the depth and duration of its convection differ. The repairs of
constants and units make these conditions enter the calculation consistently. Predicting the resulting climate
also requires a radiation scheme and cloud relations appropriate to those conditions. The measured
sensitivities identify where that further work could change the climate most in the configurations tested.

Radiation offers a way to extend the calculation through the properties of the gases themselves. The in-house
line-by-line model has already supplied the reference for the lunar calibrations; a reduced scheme built from
its spectroscopy would compute absorption and cooling over the columns for which it was checked. An independent
comparison of that spectroscopy, followed by checks of the reduced scheme by layer, would define its usable
pressures, temperatures and gases under the chosen starlight. Within that range, a modeller could change a
world's column and follow where the added gas absorbs and emits, with the vertical heating computed from the
same molecular data throughout. This would also make comparisons between worlds less dependent on a sequence
of separate broadband refits.

Cloud requires a different relation between scales. A GCM determines the large-scale circulation in which
convection develops, while a cloud-resolving run can compute the condensate and precipitation that the GCM
parameterizes. The forcing methods developed for the lunar rings and boxes make that comparison possible on
slices of a GCM climate. Their use could be extended to calibrating the GCM's cloud against resolved runs of
the same world, so that a change in rotation or gravity would be followed through the convection into the cloud
that heats and cools the global atmosphere. Because the resolved runs retain their own empirical microphysics,
this comparison also moves the question of accuracy down to the particles from which the cloud and rain form.

Rain at weak gravity makes that dependence tangible. The fall-speed repair slows drops according to the balance
of weight and drag, but CM1 still bounds their sizes and breakup with Earth's values. In the lunar box most of
the water leaving the clouds evaporated before reaching the ground, and the surface received about a quarter
of the forcing GCM's rain. That disagreement joins two possible causes: the GCM's parameterized convection and
the resolved model's treatment of falling water. Larger drops could survive more of the descent, while the
capillary scaling of the largest stable drop leaves open the sizes produced by collisions. Testing those size
limits would determine how much of the rain deficit comes from the retained Earth settings, before the resolved
rain is used to revise the lakes' water budgets.

The lake question connects a planetary design to physics that can be investigated at the scale of individual
particles. A similar connection is already explicit in the charging calculations. Two laboratory charging laws
produced different arrangements of charge in windows of the lunar storm runs because they disagree on the sign
rimed ice takes in little cloud water. Measurements at the simulated collision speeds, temperatures and water
contents could discriminate between those laws for the relevant particles. The measurements would constrain
the storms' electrical structure; their ground strikes would still require a model of the leader's descent
through the deep air. Even for a world that exists as a design, some of the distinctions between its possible
weather can therefore be investigated in an Earth laboratory.

Such process tests complement the checks available within the codes. The wave-model similarity runs, for
example, test whether changing gravity produces the corresponding change of lengths and times along the
tested paths. They retain the empirical growth and breaking coefficients, whose adequacy under the new
conditions is a further physical question. Keeping these questions separate makes progress on an unobserved
world cumulative: arithmetic can be checked against the implemented laws, approximations compared with
independent calculations, and empirical relations tested where measurements are possible. Agreement at one
step specifies which uncertainty to investigate at the next.

For a constructed world, those uncertainties concern conditions that the designer can change. The lunar
shield's dimming was chosen to obtain a target climate, and corrections to radiative geometry and gas
absorption changed the dimming required to maintain it. A revised cloud calibration or a consistent longwave
fit could change that choice again. The calculation therefore becomes an iteration between physical
prediction and design: revised heating changes the proposed illumination, which changes the climate from
which the storm and sea models take their forcing. A design temperature becomes useful only together with
the rain, circulation and cloud produced under the illumination that maintains it.

This dependence also gives the disagreement between models a practical use. Once the missing transport across
the lunar ring was imposed, CM1 and the GCM still computed different temperatures near the ground. Comparing
their radiative cooling against the reference separated part of that difference from the remaining boundary-layer
and forcing questions. Following a disagreement through its physical terms can identify which choice in the
design depends on the disputed process. The next calculation can then test whether a change in radiation
alters the required shielding, or whether a change in drop size alters the water supplied to a lake. The
design develops through questions whose outcomes would change it.

The same method of investigation can extend to natural planets, where gravity, rotation and starlight are
conditions to infer or prescribe. The repaired codes have so far simulated one world away from Earth's
values, apart from numerical controls, so their wider use begins with further tests. Yet the lunar case has
already identified concrete connections to examine: rotation changes the distribution of heating and the
cloud that responds to it; gravity changes both the depth of the atmosphere and the rain falling through it;
the height of the air relative to the radius changes which layers remain illuminated after the ground enters
shadow. Varying these properties across worlds would test how the processes combine, including regimes in
which an Earth calibration succeeds for a different reason or ceases to describe the controlling process.

The prospect is a comparison of worlds at the scales on which their environments take shape. A climate
calculation would specify the circulation and forcing from which storms develop; resolved storms would test
the cloud and rain assumed in that climate; winds and wave spectra would connect the atmosphere to a
particular coast. As these connections are checked, the question of whether a world is temperate can be
developed into questions about where its water falls, how its storms evolve through a long day, and what its
seas do under the resulting winds. Worlds with similar mean temperatures could then be compared through the
different physical environments those temperatures accompany.

For worlds still being designed, that comparison would precede choices about the environments to create.
Familiar air and water could be considered under a different gravity and Sun, with the consequences followed
far enough to determine whether rain could replenish a basin or swell would reach a proposed shore. The
breathable Moon already gives those questions physical form: tall storms above slow-falling rain, long-period
waves under light winds, and upper air still illuminated deep into the ground's night. Developing and testing
the models would turn such possibilities into alternatives that can be examined together, so that a world
chosen for its warmth could also be chosen with knowledge of the weather, water and sky beneath which life
would have to grow.
